\subsection{INTEGRAL CENTRAL FILTRATIONS}

Recall (no. 1, Remark) that a filtration \((\mathrm{G}_{\alpha})\) on the group G is called integral if \(\mathrm{G}_{\alpha} = \mathrm{G}_{n}\) for every integer \(n\) and all \(\alpha \in ]n - 1,n]\) . To be given an integral central filtration on a group G is equivalent to being given a sequence \((\mathrm{G}_n)_{n\geqslant 1}\) of subgroups of G satisfying the conditions

\begin{enumerate}
\def\labelenumi{(\roman{enumi})}
\tightlist
\item
\end{enumerate}

\[
\mathrm{G} _ {1} = \mathrm{G}\tag{i}
\]

\[
\mathrm{G} _ {n} \supset \mathrm{G} _ {n + 1} \quad \text { for   all } n \geqslant 1\tag{ii}
\]

\[
\left(\mathrm{G} _ {m}, \mathrm{G} _ {n}\right) \subset \mathrm{G} _ {m + n} \quad \text { for } m \geqslant 1 \text { and } n \geqslant 1.\tag{iii}
\]

For every integer \(n \geqslant 1\) , \(G_n\) is a normal subgroup of \(G\) and the quotient \(\mathrm{gr}_n(G) = G_n / G_{n+1}\) is commutative. On taking quotients, the mapping \((x,y) \mapsto (x,y) = x^{-1}y^{-1}xy\) of \(G_m \times G_n\) into \(G_{m+n}\) allows us to define on \(\mathrm{gr}(G) = \bigoplus_{n \geqslant 1} \mathrm{gr}_n(G)\) a graded Lie algebra structure of type \(N\) over the ring \(Z\) .

Recall (Algebra, Chapter I, §6, no. 3, Definition 5) that the lower central series of the group \(\mathbf{G}\) is defined by

\[
\mathrm{C} ^ {1} \mathrm{G} = \mathrm{G}, \quad \mathrm{C} ^ {n + 1} = (\mathrm{G}, \mathrm{C} ^ {n} \mathrm{G}) \quad \text { for } n \geqslant 1.\tag{28}
\]

The corresponding filtration is called the lower central filtration of G.

PROPOSITION 4. (i) The lower central series of G is an integral central filtration on G.

\begin{enumerate}
\def\labelenumi{(\roman{enumi})}
\setcounter{enumi}{1}
\tightlist
\item
  If \((\mathrm{G}_n)_{n\in \mathbf{N}^*}\) is an integral central filtration on \(\mathbf{G}\) , then \(\mathbf{C}^n\mathbf{G}\subset \mathbf{G}_n\) for all \(n\in \mathbf{N}^*\) .
\end{enumerate}

Assertion (i) has been proved in Algebra, Chapter I, § 6, no. 3, formula (7).

We prove (ii) by induction on \(n\) ; \(\mathrm{C}^1\mathrm{G} = \mathrm{G} = \mathrm{G}_1\) ; for \(n > 1\) ,

\[
\mathrm{C} ^ {n} \mathrm{G} = (\mathrm{G}, \mathrm{C} ^ {n - 1} \mathrm{G}) \subset (\mathrm{G}, \mathrm{G} _ {n - 1}) \subset \mathrm{G} _ {n}.
\]

PROPOSITION 5. Let \(\mathbf{G}\) be a group and \(\operatorname{gr}(\mathbf{G})\) the graded Lie \(\mathbf{Z}\) -algebra associated with the lower central filtration on \(\mathbf{G}\) . Then \(\operatorname{gr}(\mathbf{G})\) is generated by \(\operatorname{gr}_1(\mathbf{G}) = \mathbf{G}/(\mathbf{G},\mathbf{G})\) .

Let L be the Lie subalgebra of \(\mathrm{gr}(\mathrm{G})\) generated by \(\mathrm{gr}_{1}(\mathrm{G})\) ; we show that \(\mathrm{L} \supset \mathrm{gr}_{n}(\mathrm{G})\) by induction on n, the assertion being trivial for n = 1. Suppose that n > 1 and \(\mathrm{L} \supset \mathrm{gr}_{n-1}(\mathrm{G})\) . As \(\mathrm{C}^{n}\mathrm{G} = (\mathrm{G}, \mathrm{C}^{n-1}\mathrm{G})\) , the construction of the Lie algebra law on \(\mathrm{gr}(\mathrm{G})\) shows immediately that

\[
\operatorname{gr} _ {n} (\mathrm{G}) = \left[ \operatorname{gr} _ {1} (\mathrm{G}), \operatorname{gr} _ {n - 1} (\mathrm{G}) \right] \subset \mathrm{L}.
\]

The above proof shows that the lower central series of the Lie algebra \(\mathbf{gr}(\mathbf{G})\) (§ 2, no. 7) is given by

\[
\mathcal {C} ^ {n} (\mathrm{gr} (G)) = \sum_ {m \geqslant n} \mathrm{gr} _ {m} (G).\tag{29}
\]

Remark. Let \(k\) be a ring, \(n\) an integer \(>0\) and A set of lower triangular matrices with \(n\) rows and \(n\) columns and elements in \(k\) . For \(p \geqslant 0\) , let \(\mathbf{A}_p\) be the set of \((x_{ij}) \in \mathbb{A}\) such that \(x_{ij} = 0\) for \(i - j < p\) . Then \(\mathrm{A}_0 = \mathrm{A}\) and \(\mathrm{A}_p\mathrm{A}_q \subset \mathrm{A}_{p + q}\) . Let \(\Gamma_p = 1 + \mathrm{A}_p\) . Then \(\Gamma_1\) is a subgroup of \(\mathbf{GL}(n,k)\) called the strict lower triangular group of order \(n\) over \(k\) . By Proposition 2 of no. 5, \((\Gamma_p)\) is an integral filtration on \(\Gamma_1\) . As \(\Gamma_n = \{1\}\) , we see that the group \(\Gamma_1\) is nilpotent (Algebra, Chapter I, §6, no. 3, Definition 6).

